\section{Local invariant rings and balanced representatives}
\label{app:local-algebra}

This appendix records two local calculations used in
Proposition~\ref{prop:physical-kernel}.

\subsection[E3 maximal-torus invariants]{$E_3$ maximal-torus invariants}

The closure of the minimal nilpotent orbit of $\mathfrak{sl}_n$ has the
rank-one presentation
\begin{equation}
 \overline{\mathcal O_{\min}(\mathfrak{sl}_n)}
 =\Spec\left(
 \frac{\CC[q_1,\ldots,q_n,\widetilde q_1,\ldots,\widetilde q_n]}
      {(\sum_iq_i\widetilde q_i)}
 \right)^{\CC^*_{\mathrm{ADHM}}},
 \label{eq:app-adhm}
\end{equation}
where the ADHM torus acts with weights $+1$ and $-1$ on $q_i$ and
$\widetilde q_i$. The complex moment-map matrix is
\begin{equation}
 \Phi^i{}_j=q_i\widetilde q_j,
 \qquad \tr\Phi=0.
\end{equation}
The flavor torus acts by
$q_i\mapsto t_iq_i$ and
$\widetilde q_i\mapsto t_i^{-1}\widetilde q_i$, with $\prod_it_i=1$.
After combining it with the ADHM torus, an invariant monomial has equal powers
of $q_i$ and $\widetilde q_i$ for each $i$. The invariant ring is therefore
generated by
\begin{equation}
 z_i=q_i\widetilde q_i,
 \qquad \sum_i z_i=0.
 \label{eq:app-zi}
\end{equation}
Since a complex torus is linearly reductive, taking invariants is exact and
introduces no additional relation. Thus
\begin{align}
 \CC[\overline{\mathcal O_{\min}(\mathfrak{sl}_3)}]^{T_{SU(3),\CC}}
 &\cong\frac{\CC[z_1,z_2,z_3]}{(z_1+z_2+z_3)}
 \cong\CC[H_1,H_2],\\
 \CC[\overline{\mathcal O_{\min}(\mathfrak{sl}_2)}]^{T_{SU(2),\CC}}
 &\cong\frac{\CC[w_1,w_2]}{(w_1+w_2)}
 \cong\CC[L].
\end{align}
One convenient root-basis convention is
\begin{equation}
 H_1=z_1,
 \qquad H_2=-z_3,
 \qquad L=w_1=-w_2.
 \label{eq:app-root-coordinates}
\end{equation}
The $E_3$ chiral ring is the fiber product of the two orbit rings over their
origins. Its invariant ring is consequently
\begin{equation}
 \CC[H_1,H_2]\mathop{\times}_{\CC}\CC[L]
 \cong\frac{\CC[H_1,H_2,L]}{(H_1L,H_2L)}.
\end{equation}

\subsection{Balanced real moment maps}

For arbitrary $(H_1,H_2)\in\CC^2$, solve
\eqref{eq:app-zi} and \eqref{eq:app-root-coordinates} by
\begin{equation}
 z_1=H_1,
 \qquad z_2=-H_1+H_2,
 \qquad z_3=-H_2.
\end{equation}
For each $z_i=|z_i|e^{\ii\theta_i}$ choose
\begin{equation}
 q_i=\sqrt{|z_i|}e^{\ii\theta_i},
 \qquad
 \widetilde q_i=\sqrt{|z_i|}.
 \label{eq:app-balanced-pairs}
\end{equation}
Then $q_i\widetilde q_i=z_i$ and
$|q_i|^2-|\widetilde q_i|^2=0$ for every $i$. The ADHM real moment map and all
flavor-Cartan real moment maps vanish. The $SU(2)$ component is identical with
$w_1=L,w_2=-L$. Equations \eqref{eq:app-balanced-pairs},
\eqref{eq:conifold-balanced} and \eqref{eq:e2-moment-maps} supply the local
balanced lifts used in the proof of Proposition~\ref{prop:physical-kernel}.

\subsection[Abelian monopole sums for the E3 components]{Abelian monopole sums for the $E_3$ components}

Removing the diagonal magnetic charge from the affine-$A_1$ quiver leaves
$m\in\ZZ$. Its monopole sum is
\begin{equation}
 \frac1{1-x}\sum_{m\in\ZZ}x^{|m|}
 =\frac{1+x}{(1-x)^2}.
\end{equation}
For the affine-$A_2$ triangle, set one diagonal charge to zero. The sum is
\begin{equation}
 \frac1{(1-x)^2}
 \sum_{m_1,m_2\in\ZZ}
 x^{\frac12(|m_1-m_2|+|m_1|+|m_2|)}
 =\frac{1+4x+x^2}{(1-x)^4}.
\end{equation}
The two series are respectively the character dimensions of the
$SU(2)$ representations of highest weight $2n$ and the $SU(3)$
representations with Dynkin labels $(n,n)$. Adding the two component rings and
subtracting their common constant gives \eqref{eq:e3-hs-union}.
